\documentclass[11pt]{article}
\input{packages.tex}
\usepackage{ulem}
\DeclareMathOperator{\sign}{sign}
\DeclareMathOperator{\Imm}{Im}
\DeclareMathOperator{\DIF}{DIF}
\DeclareMathOperator{\rg}{rg}
\DeclareMathOperator{\dist}{dist}

% Настройка колонтитулов
\pagestyle{fancy}
\fancyhf{} % Очистить все поля
\fancyhead[L]{Пустяки, дело-то житейское:) }
\fancyhead[R]{\href{ https://t.me/mallina_olga  }{Автор}}
\fancyfoot[C]{\thepage}               % Справа внизу

\setlength{\headheight}{13.6pt}

\hypersetup{
    colorlinks=true,
    linkcolor=blue,
    citecolor=blue,
    filecolor=blue,
    urlcolor=blue,
}

\begin{document}


\begin{titlepage}
    \centering
    \vspace*{\fill}  % заполняет пространство сверху

    {\Huge Лекция 9. Машины Тьюринга. Вычислимость и разрешимость. 
    \par}

    \vspace*{\fill}  % заполняет пространство снизу
\end{titlepage}

\setcounter{page}{1}

\section{Пример машины Тьюринга}
      \begin{figure}[H]
        \centering
        \includegraphics[width=0.7\textwidth]{pictures/пример машины тьюринга.png}
        
    \end{figure}

\section{Конфигурации машины Тьюринга}



\begin{definition}
    Расширенной конфигурацией машины Тьюринга назовём двустороннюю бесконечную последовательность $\ldots \Lambda w \Lambda \ldots,$ где
    $w$ --- конфигурация машины Тьюринга, а на месте многоточий стоят пустые символы $\Lambda$. \\
    Отображение на расширенных конфигурациях МТ:
    \[ d_M (uqv) = 
    \begin{cases}
        ua'q'v', \quad \delta(a,q) = (a',+1,q')\\
        uq'a'v', \quad \delta(a,q) = (a',0,q')\\
        u'q'ba'v', \quad \delta(a,q) = (a',-1,q') 
    \end{cases} 
    \]
    где $v=av', \ u=u'b$
\end{definition}

      \begin{figure}[H]
        \centering
        \includegraphics[width=0.6\textwidth]{pictures/конфигурация мт.png}
        
    \end{figure}
\begin{figure}[H]
    \centering
    \includegraphics[width=0.6\textwidth]{pictures/конфигурация мт2.png}
        
\end{figure}

\begin{definition}
    Отображение на конфигурациях формализует работу МТ за один такт. Работа машины Тьюринга М
    полностью задаётся начальной конфигурацией и отображением на конфигурациях и может быть описана конечной или бесконечной последовательностью конфигураций.
    \[ w_0, w_1,..., w_n, ... \text{, где } w_{i+1}=d_M(w_i) \]
    Последовательность $w_0, w_1, ..., w_n, ...$ называется $\textbf{протоколом работы}$ МТ.
\end{definition}

\begin{definition}
    $\textbf{Результатом работы }$ МТ является слово, которое получается из последней конфигурации, вычёркиванием символа состояния управляющего устройства. 
\end{definition}

\begin{definition}
    Частично определённая функция 
    \[ f: \mathbb{N}^n \to \mathbb{N} \]
    называется вычислимой на машине Тьюринга, если существует такая машина Тьюринга М, что для любого набора $(m_1,...,m_n)$, принадлежащего области определения $f$, 
    работа $M$ на входе 
    \[ \#\underbrace{11\ldots1\#}_{m_1+1}\,
    \underbrace{11\ldots1\#}_{m_2+1}\,\ldots\,
    \#\underbrace{11\ldots1\#}_{m_n+1} \]
    Завершается с результатом
    \[ \underbrace{11...1}_{f(m_1,m_2,...,m_n)+1} \]
\end{definition}

\section{Вычислимость и разрешимость}
\begin{definition}
    $\textbf{Решающая}$ машина Тьюринга имеет ровно два финальных состояния: $q_{yes}$ и $q_{no}$.\\
    Язык $L$ в конечном алфавите $A$ называется $\textbf{разрешимым}$, если существует такая решающая МТ с алфавитом $B\supset A$, которая на любом входе $w\in L$ останавливается в состоянии $q_{yes}$, а на любом входе $w\notin L$ останавливается в состоянии $q_{no}$. 
    При этом пустой символ $\Lambda$ не принадлежит алфавиту $A$.
\end{definition}


\begin{definition}
    Характеристическая функция $ \chi_L$ языка $L$:
    \[ \chi_L = 
    \begin{cases}
        1, w\in L\\
        0, w\notin L
    \end{cases} \]
\end{definition}

\begin{remark}
    Разрешимость языка равносильна вычислимости его характеристической функции.
\end{remark}



\section{Кодирование и моделирование}

МТ $M'$ моделирует работу М, если она удовлетворяет следующему свойству: на входе вида $c(w)$, где $w$ --- вход М,
результатом работы $M'$ является $c(u)$, где $u$ --- результат работы М на входе $w$.\\
Здесь $c(x)$ --- код слова x, то есть отображение из множеста слов над алфавитом М в множество слов над алфавитом $M'$, позволяющее по $c(x)$ однозначно восстановить х.
\[ w\xrightarrow{M} u \]
\[ c(w)\xrightarrow{M'} c(u) \]

\begin{theorem}
    Для любой машины Тьюринга $M$ существует машина в двухбуквенном алфавите $M'$, моделирующая работу М.
\end{theorem}

\begin{theorem}
    Для любой k-ленточной машины Тьюринга M существует одноленточная машина Тьюринга $M'$, моделирующая работу M. 
\end{theorem}

\begin{definition}
    $\textbf{Универсальной}$ машиной Тьюринга называют такую МТ, которая моделирует работу любой другой МТ.\\
    Входом такой универсальной машины Тьюринга является описание машины Тьюринга и её входа.
\end{definition}

\begin{theorem}
    Существует универсальная машина Тьюринга с одной лентой и двухбуквенным алфавитом.
\end{theorem}


\section{Алгоритмически неразрешимые проблемы}
\begin{itemize}
    \item Проблема самоприменимости 
    \item Проблема остановки
    \item Проблема достижимости в ассоциативном исчислении
    \item Проблема равенства слов в полугруппе
    \item Проблема соответствия Поста
    \item Функция трудолюбия Радо
    \item Неразрешимость множества общезначимых формул
\end{itemize}

\subsection{Проблема самоприменимости}
      \begin{figure}[H]
        \centering
        \includegraphics[width=0.7\textwidth]{pictures/проблема самоприменимости.png}
        
    \end{figure}
Множество алгоритмов счётно. Зафиксируем их нумерацию $A_0,A_1,...$ (повторения допустимы).\\ 
Число входов тоже счётно, потому что без ограничения общности считаем, что на вход алгоритму подаётся натуральное число. \\
Определим функцию $ F: \mathbb{N}\to \mathbb{N} $ таким образом: $F(i)=1$, если алгоритм $A_i$ останавливается на входе $i$, и $F(i)=0$ в противном случае.\\
Проблема самоприменимости состоит в вычислении функции $F$. Она алгоритмически неразрешима при любой нумерации алгоритмов. 

\begin{proof}
    От противного: пусть $F(i)$ --- вычислима $\Rightarrow$ есть машина Тьюринга М, которая её вычисляет. 
    Построим $M'$ : если $M(i)=0$, то $M'(i)=1$, а если $M(i)=1$, то $M'$ на входе $i$ не останавливается вообще. 
    Но у $M'$ есть номер m. \\
    Рассмотрим $F(m)$. Если $M'$ останавливается на входе m(то есть на своём номере), то по определению функции $F \ F(m)=1$. Но тогда по построению $M'$ не останавливается на m $\Rightarrow$ получаем противоречие.\\
    Если $M'$ не останавливается на m $\Rightarrow$ по определению $F \ F(m)=0$, но тогда по построению $M' \ M'(m)=1$ (то есть останавливается)$\Rightarrow$ получаем противоречие $\Rightarrow$ функция $F$ невычислима.
\end{proof}


\subsection{Следствия: проблема остановки}
Определим функцию $f_s: \mathbb{N}\to \mathbb{N}$ таким образом: $f_s(i,j) =1$, если алгоритм $A_i$ останавливается на входе $j$, и $f_s(i,j)=0$ в противном случае.\\
Проблема остановки заключается в вычислении функции $f_s$. Она алгоритмически не разрешима при любой нумерации алгоритмов.\\
Обозначим через stop язык в алфавите $\{0,1\}$, состоящий из тех описаний $\langle M, w \rangle$ машины Тьюринга M и входного слова $w$, для которых M останавливается на входе $w$. \\
Язык stop неразрешим.


\section{m-сводимость языков}
Рассмотрим два языка $A$ и $B$. Задача заключается в разрешении этих языков. 
Будем говорить, что язык $A\in \mathbb{N}$ m-сводится к языку $B\in \mathbb{N}$  ($A\le_m B$), если существует всюду определённая вычислимая функция $f: \mathbb{N}\to \mathbb{N}$ такая что, 
\[ x\in A \Leftrightarrow f(x)\in B \]



\end{document}


